% =============================================================================
% 4 SEQUENTIAL VS. AGILE PLANNING            budget: 1.25 pages - the SHORTEST
%                                            main chapter
% Answers TASK QUESTION 3, verbatim:
%   "Compare how the further planning of the project would take place in a
%    sequential approach and an agile approach. In doing so, refer to the
%    specific project and do not describe the procedure in general terms."
%
% -----------------------------------------------------------------------------
% LITERATURE: MIXED, and the best place in the report to show OWN RESEARCH
%   The structure to aim for in every paragraph: one cited sentence describing
%   how the two paradigms differ on this axis, then two or three uncited
%   sentences playing it out on a named recipe-app feature. The citation
%   carries the claim, the example carries the marks.
%   Cite here:
%     - the SEQUENTIAL side, which is the half you have not written about
%       anywhere else: PMBOK 7 and Kerzner (2009) are the named sources -
%       "you can find more about sequential projects or waterfall projects in
%       the project management body of knowledge" [G-Jun26]. A tutor pointed
%       explicitly at Kerzner 12th ed., section 8.12 "Agile Project
%       Management", p. 287 - CAUTION: the task PDF lists Kerzner (2009), the
%       10th edition. If that locator is used, cite the 12th edition as its own
%       entry; do not hang p. 287 on the 2009 edition
%     - Wysocki, "Effective Project Management: Traditional, Agile, Extreme,
%       Hybrid" (8th ed., 2019) - named by a tutor among the extra literature,
%       and the obvious fit for this comparison; it also underwrites the hybrid
%       ending of chapter 5
%     - the iron triangle and its inversion - find a proper source for the
%       constraint model rather than asserting it; this is the sharpest
%       argument of the chapter and should not stand uncited
%     - the Scrum events you describe here (review, retrospective, daily):
%       Scrum Guide, one citation covering all three
%   Research, best case - this is where reading is visibly rewarded:
%     - one or two PEER-REVIEWED comparisons of sequential and agile
%       development, ideally empirical rather than opinion. The
%       "sequential_vs_agile__*" and "evidence__*" PDFs in literature/ are
%       exactly this. Two solid papers here do more for the grade than ten
%       textbook citations elsewhere
%     - anything measuring requirements volatility or the cost of late change
%       supports axes 1, 2 and 5 at once
%   NOT cited: which recipe feature would be cut before the deadline, how your
%   fictitious project would have run under waterfall.
%   Do NOT: cite consultancy websites or "waterfall vs agile" blog posts - a
%   commercial web basis is explicitly warned against [G-Jun26].
%
% WHAT YOU WANT TO SAY HERE, IN SHORT
%   - "Under a sequential approach this project would have frozen its
%      requirements in a specification; under agile the backlog keeps
%      changing - and here is what that means for THIS feature." [cite + own]
%   - "The sequential plan claims to know the whole project up front; the agile
%      plan only claims accuracy for the next few weeks." [cite + own]
%   - "There, the manager estimates; here, the people doing the work do."
%      [cite + own]
%   - "There, one delivery at the end; here, a usable increment every sprint -
%      and users see it at the sprint review." [cite the events + own]
%   - "The constraints invert: sequential fixes scope and lets time and cost
%      move, agile fixes time and cost and lets scope move. Concretely: this
%      feature is the one that gets dropped." [cite + own - strongest point]
%   - "Agile is the right choice for this product because requirements will
%      change and time to market decides; for a safety-critical system it
%      would not be." [own, one source if available]
% -----------------------------------------------------------------------------
% THE MOST-REPEATED WARNING OF THE ENTIRE MODULE - DO NOT GO GENERIC
%   "Don't just put in a generic comparison" [G-Mar25]
%   "Describe specifically how you would have done THIS project if you had
%    followed a sequential approach" [G-Sep25]
%   Repeated in [G-Nov25], [G-Dec25], [G-Jan26] and [S4], and restated in the
%   task PDF itself. A Schlömer-era ruling that allowed a "more general"
%   comparison [SCH, 20.03.23] belonged to the OLD task set and is void.
%   => Every axis below must be played through on a NAMED recipe-app feature
%      from Table~\ref{tab:user-stories}. The axes are free - "you've got
%      plenty of leeway in how you compare and contrast" [G-Jan26] - the
%      concreteness is not.
%
% KEEP IT SHORT: "this third sub sub task can definitely be shorter, so focus
% more on the first and the second" [S4]. It is an OUTLOOK on how the sprints
% after the first one would be planned. Do NOT build a second complete
% sequential project plan [S1] - no Gantt chart, no phase schedule, no WBS.
%
% THIS IS THE ONLY PLACE to describe the Scrum ceremonies in detail [G-Dec25]:
% sprint review, retrospective, daily scrum. Refinement and sprint planning
% were covered in chapter 3 - do not describe them twice; refer back instead.
% -----------------------------------------------------------------------------
% AXES - pick 4 or 5, each applied to a named feature:
%   1. REQUIREMENTS: months of analysis ending in a frozen functional
%      specification, which "is out of date within a year or two years"
%      [G-Jan26], versus a backlog that keeps changing. Notation contrast:
%      use cases and UML versus user stories [G-Dec25].
%   2. PLANNING HORIZON: everything planned up front versus a plan that is
%      "only accurate for the first maybe three months" [G-Nov25].
%   3. ESTIMATION: the project manager estimates end-to-end versus the
%      developers who will do the work estimating it themselves.
%   4. DELIVERY CADENCE: one deliverable after up to two years versus working
%      software every few weeks, each increment usable - "not a prototype that
%      you throw away".
%   5. THE INVERTED IRON TRIANGLE [S2] - the sharpest single argument:
%      sequential fixes SCOPE and lets time and cost vary; agile fixes TIME and
%      COST and lets scope vary. Apply it concretely: as the deadline nears, a
%      minor recipe feature is dropped while the core ships.
%   6. FEEDBACK AND RISK: sequential verifies only at the end [S2]; stakeholders
%      are "not actively involved, so we don't get any feedback" [S3]. The
%      sprint review is open to everyone, so "at the end of the development,
%      the solution is not something that they've never seen before" [G-Feb25].
%      IMPORTANT DISTINCTION: the retrospective is NOT for users - team and PO
%      only [G-Feb25].
%   7. COMMUNICATION: agile makes everything public and transparent; in
%      waterfall "there was scope for miscommunication, misinterpretation, and
%      not that much scope for active feedback" [G-Feb25].
%
% WHY AGILE FITS THIS PROJECT - close the chapter with this:
%   consumer apps "all change"; speed to market and first-mover advantage
%   dominate. Pre-internet "the cost and the time were less important than the
%   quality", post-internet "there was only one thing that mattered... how
%   quickly can you get it out into the market" [G-Jan26].
%   COUNTER-EXAMPLE he uses himself: an AIR TRAFFIC CONTROL system - fixed
%   requirements, safety-critical, regulated - is where sequential wins.
%   Naming a case where the chosen method would be wrong is a Transfer point
%   and prepares the balanced ending of chapter 5.
%
% TERMINOLOGY - he corrects these:
%   - "Waterfall" and "sequential" are synonyms [G-Dec25].
%   - "Systems engineering" and "SDLC" are UMBRELLA terms covering BOTH
%     paradigms, not a third method and not synonyms for waterfall [G-Feb25].
%   - [S1] vocabulary may be used SPARINGLY as vocabulary only: phase gates,
%     V-model, WBS, Gantt, and the classic role set (client, steering
%     committee, project manager, team, project office, suppliers).
%
% SOURCES: PMBOK is the natural source for the sequential side - "you can find
% more about sequential projects or waterfall projects in the project
% management body of knowledge" [G-Jun26]. Kerzner (2009) and Wysocki,
% "Effective Project Management: Traditional, Agile, Extreme, Hybrid" (8th ed.,
% 2019) were both named by tutors for exactly this chapter; Kerzner 12th ed.
% section 8.12 "Agile Project Management", p. 287, was pointed at explicitly.
% Citing an IU course book from another module is tolerated - "Not generally
% cited, but I think it is ok! It's better if you can find an original source"
% [G, 18.04.25].
%
% DO NOT: draw conclusions here - the overall reflection belongs to chapter 5,
% which must not duplicate this comparison.
% =============================================================================

% #############################################################################
% WRITING PLAN - CHAPTER 4
% Budget: 1.25 pages = ~875 WORDS across 7 blocks (opening + 5 axes +
% closing) => ~125 WORDS EACH. About 12-13 citations.
%
% ** THE FORMAT IS FIXED, AND IT IS WHAT MAKES THIS CHAPTER WORK: **
%   per axis - ONE cited sentence on how the two paradigms differ,
%              then 2-3 UNCITED sentences playing it out on a NAMED recipe-app
%              feature from Table~\ref{tab:user-stories}.
%   The citation carries the claim; the example carries the marks.
%
% ***** DECIDED: WRITE AXES 1, 2, 3, 5 AND 6. DELETE THE STUBS FOR 4 AND 7. *****
% Axes 1, 3, 5 and 6 are the core; 5 is the sharpest single argument in the
% report. Axis 2 is the chosen fifth because the task sheet's own background
% paragraph is about exactly this - "The planning horizon shifts so that only
% the foreseeable future is planned in detail and planning is therefore
% iterative rather than a one-time process." Answering the assignment's own
% framing is worth more than the delivery-cadence or communication axis.
%
% THE NAMED FEATURE FOR EACH AXIS IS ALREADY FIXED - see PROJECT_FACTS.md, s. 9:
%   axis 1 -> US-06   axis 2 -> US-08   axis 3 -> US-07
%   axis 5 -> US-09 (the feature that gets dropped)   axis 6 -> US-03
%
% Sources for the SEQUENTIAL half - this is the half with no coverage anywhere
% else in the report: kerzner2022 (5 of the 7 axes), pmbok2021, pmi2017.
% Sources for the AGILE half: scrumguide2020, kelly2019, beck2001*.
% Peer-reviewed comparison and counter-evidence: mishra2023.
%
% -----------------------------------------------------------------------------
% !! FOUR SOURCE-LEVEL TRAPS IN THIS CHAPTER - all verified:
%
% 1. mishra2023 ARGUES THE OPPOSITE OF AXIS 5, twice. Its Table 4 (p. 1515) and
%    p. 1517 both say AGILE's cost is the adaptable one and that waterfall suits
%    limited budgets. NEVER cite it near the inverted triangle, and never put
%    its 42.62 % cost-overrun figure in the same chapter as the inversion.
% 2. kerzner2022 Table 2-13 (p. 86) is "Traditional versus LEAN", not versus
%    agile. Use Table 8-9 (p. 307), which is explicitly the agile comparison.
% 3. The whole indented block on kerzner2022 p. 308 is Charles Cobb (2015,
%    p. 115), per footnote 7 - four paragraphs, not three lines. Kerzner's own
%    adjacent sentences are safe and are the ones used below.
% 4. mishra2023 must NEVER be used for anything Scrum: it lists five Scrum roles
%    and has "weekly sprint reviews" inside multi-week sprints.
%
% !! EDITION FOOTNOTE: add one line noting that the task sheet lists Kerzner
%    (2009, 10th ed.) while the edition read and cited here is the 13th (2022).
%    Honest, costless, and it protects the page-number criterion.
%
% !! MANIFESTO FOOTNOTE: the first Beck citation needs the locator footnote -
%    see the wording in 05_reflection_and_conclusion.tex.
% #############################################################################

\section{Sequenzielle vs. agile Planung der Rezept-App}

% --- OPENING PARAGRAPH: framing and terminology --------------------------
% ONE CITATION. Two or three sentences, no more.
%
%   \parencite[PMBOK Guide, Section~2.3.3]{pmbok2021}
%   "A predictive approach is useful when the project and product requirements
%    can be defined, collected, and analyzed at the start of the project. This
%    may also be referred to as a waterfall approach."
%   -> Does two jobs: it establishes that "sequential", "predictive" and
%      "waterfall" name the same thing (the task says "sequential", most sources
%      say "waterfall"), and it states the sequential precondition neutrally
%      before the chapter argues against it.
%
% OWN VOICE: say what this chapter does - it compares how the planning of THIS
% project would have continued under each approach. It is an OUTLOOK on the
% sprints after the first one, not a second complete project plan. No Gantt
% chart, no phase schedule, no WBS.
Kapitel~3 führt die Planung von Cookbox bis zum Ende des ersten Sprints; dieses
Kapitel vergleicht, wie sie unter einem sequenziellen und unter einem agilen
Vorgehen weiterliefe. Was das Project Management Institute (PMI) einen predictive
approach nennt, \enquote{is useful when the project and product requirements can be
defined, collected, and analyzed at the start of the project} und kann
\enquote{also be referred to as a waterfall
approach} \parencite[PMBOK Guide, Section~2.3.3]{pmbok2021}. Es folgt ein
Ausblick auf die Sprints nach dem ersten und kein zweiter Projektplan, und jede
der fünf nachstehenden Achsen wird an einer Story aus
Tabelle~\ref{tab:user-stories} durchgespielt.

Traditionelles Projektmanagement ist vorschreibend, mit klar definiertem Scope und
sequenziellen Lebenszyklusphasen, und \enquote{The intent is to prevent changes in
the deliverables} \parencite[p.~308]{kerzner2022}.\footnote{Die Aufgabenstellung
nennt Kerzner (2009, 10. Aufl.); gelesen und zitiert wird hier die 13. Auflage
(2022).}
Die agile Position ist die entgegengesetzte, und sie wird auf beiden Ebenen
formuliert: Das Manifest
schätzt \enquote{Responding to change over following a plan}
\parencite[para.~2]{beck2001manifesto}, was das zweite Prinzip zu
\enquote{Welcome changing requirements, even late in development} ausführt
\parencite[para.~2]{beck2001principles}.\footnote{Das Agile Manifest und die
zugehörigen Prinzipien sind ausschließlich als Webseiten ohne Seitenzählung
veröffentlicht, weshalb APA~7 eine Zitation nach Absatz statt nach Seite
verlangt; die vier Werte und zwölf Prinzipien erscheinen in der Quelle
unnummeriert, und die hier verwendete Zählung stammt vom Verfasser und beginnt
beim ersten Prinzip.} Unter einem sequenziellen Vorgehen wäre US-06, die
Einkaufsliste, in einer Spezifikation festgeschrieben worden, die Monate vor der
ersten Codezeile abgenommen wurde, samt ihrer Abhängigkeit von US-05. Im Backlog
wurde sie umformuliert, nachdem Folgegespräche gezeigt hatten, dass die meisten
Teilnehmer ohnehin auf Papier planen und nur die Liste wollten, und die
Abhängigkeit löste sich damit auf. Eine abgenommene Spezifikation hätte dieselbe
Information als Änderungsantrag verarbeitet.

Sequenzielle Planung beansprucht, das ganze Projekt von Anfang an zu kennen.
\textcite[p.~378]{kerzner2022} hält fest, dass eine vollständige Programmplanung
\enquote{cannot be accomplished unless all of the necessary information becomes
available at project initiation}, wozu der Statement of Work und der
Projektstrukturplan gehören, obwohl er an anderer Stelle die Planung als iterativ
und fortlaufend bezeichnet
\parencite[p.~367]{kerzner2022}. Kelly beziffert den agilen Horizont:
\enquote{prioritizing into the far future (e.g., more than 12 weeks out)}
verringert die Agilität \parencite[ch.~8, pp.~57--62]{kelly2019}. Keines dieser
Dokumente existierte bei Cookbox zu Projektbeginn, und nichts in einer
Sechsmonatsspezifikation hätte die datenschutzrechtliche Prüfung vorweggenommen,
die US-08, die Kontolöschung, vom Could-have zum Must-have machte. Der Plan war
etwa zwei Sprints lang belastbar.

Im traditionellen Projektmanagement kommen die Zahlen von Managern: Um die
verstrichene Zeit zwischen zwei Ereignissen zu bestimmen, müssen
\enquote{responsible functional managers evaluate the situation and submit their
best estimates} \parencite[p.~455]{kerzner2022}. In Scrum schätzen die Entwickler, die einen
Eintrag bauen werden, ihn selbst, und diese Umkehrung hat dieses Backlog
verändert. Ein Manager, der US-07, die Umrechnung der Portionszahl, nach ihrer
einzeiligen Beschreibung bepreist hätte, hätte sie niedrig angesetzt. Die vier
Entwickler schätzten die Einheitenumrechnung über Gramm, Milliliter, Tassen und
ganze Eier deutlich über der ursprünglichen Annahme, und die Story rückte an das
Ende der Should-haves (Abschnitt~\ref{subsec:refinement}). Die Schätzung hat die
Reihenfolge verändert und nicht umgekehrt.

Anschließend kehren sich die Randbedingungen um, und am schärfsten formuliert das
eine traditionelle Autorität. \textcite[p.~78]{kerzner2022} stellt die beiden
Fälle nebeneinander: Der Wasserfallansatz \enquote{begins with well-defined requirements from
which we must determine the budget and schedule to produce the deliverables},
während es zum anderen heißt: \enquote{we may start out with a fixed budget and schedule,
and then have to decide how much work can be done within the time and cost
constraints. The requirements can evolve over the life of the project.} Seine
Tabelle der Randbedingungskategorien benennt beide \parencite[p.~609]{kerzner2022}: Ein
sequenzielles Cookbox entspricht Situation A-3, in der die Leistung feststeht,
während Zeit und Kosten variieren; das Scrum-Projekt ist B-1, wo stattdessen
Zeit und Kosten feststehen. Kelly leitet dieselbe Umkehrung aus den Kosten als
Produkt aus Personen und Zeit ab: Personen lassen sich kurzfristig nicht
hinzufügen, die Qualität muss hoch bleiben, und der Sprint legt die Zeit fest,
also gilt: \enquote{That leaves scope} \parencite[ch.~1, pp.~5--10]{kelly2019}. Sechs Monate
und 240.000~EUR stehen hier fest, und das erste Release bringt nichts ein, also
gibt der Scope nach: Gestrichen wird US-09, das Moderationswerkzeug, ein
Could-have aus Tabelle~\ref{tab:initial-backlog}, das das Support-Team bei den
Startvolumina von Hand bewältigen kann, während die Must-haves --- US-08 nun
darunter --- zum Termin ausgeliefert werden.

\textcite[p.~307]{kerzner2022} verzeichnet das Kundenfeedback im traditionellen
Projektmanagement als \enquote{Minimal, perhaps only at project termination} und im
agilen als \enquote{Throughout the project}. Drei der Prüfpunkte, die Scrum in
jeden Sprint einbaut, berühren diesen Unterschied: ein fünfzehnminütiger Daily
Scrum für die Entwickler, ein Sprint
Review, in dem \enquote{The Scrum Team presents the results of their work to
key stakeholders} \parencite[p.~9]{scrumguide2020}, und eine Retrospektive, für
die der Guide allein das Scrum Team als Teilnehmende benennt
\parencite[p.~10]{scrumguide2020}. US-03 gewinnt am meisten: Eine Kochansicht,
die jeweils einen Schritt in großer Schrift zeigt, ist diejenige Story, die sich
aus einer schriftlichen Beschreibung am ehesten falsch bauen lässt, und ein
Stakeholder, der sie im Review auf einem Telefon in der Hand hält, lernt mehr,
als jede Spezifikation vermittelt. Unter einem sequenziellen Plan fällt dieser
erste Blick von außen erst bei der Abnahme nach fünf Monaten, und fünf Monate
sind zugleich die Arbeit, die ungeprüft in die falsche Richtung laufen kann.

PMI benennt vier Arten von Projekten, für die iterative und agile Vorgehensweisen
gut geeignet sind: solche, die \enquote{Require research and development; Have high rates of
change; Have unclear or unknown requirements, uncertainty, or risk; or Have a
final goal that is hard to describe} \parencite[p.~16]{pmi2017}. Cookbox erfüllt
die mittleren beiden klar: Drei Refinement-Entscheidungen haben allein im ersten
Sprint Band, Formulierung oder Reihenfolge einer Story verändert. Die vierte
erfüllt es nur locker und die erste gar nicht. Zu PMIs eigenen ausgearbeiteten
Fällen gehört ein militärisches Nachrichtensystem, bei dem \enquote{Many lives
were potentially at risk} und die Anforderungen über Zulieferorganisationen
hinweg festgeschrieben waren
\parencite[p.~138]{pmi2017}; Cookbox ist das entgegengesetzte Projekt. Agilität
ist zudem nicht kostenlos: \textcite[p.~1516]{mishra2023} geben in einem
narrativen Review den Befund einer anderen Studie wieder, wonach agile Projekte
ihre Terminpläne um 22,4\% überschreiten gegenüber 13,65\% beim Wasserfallmodell.


% --- Axis 1 [CORE]: frozen specification vs. changing backlog ------------
% TWO CITATIONS, one per side.
%
% SEQUENTIAL:
%   \parencite[p.~308]{kerzner2022}
%   "Traditional project management is viewed as prescriptive with well-defined
%    scope and life-cycle phases where work is performed sequentially. The
%    intent is to prevent changes in the deliverables."
%   -> The cleanest one-sentence definition available. "The intent is to prevent
%      changes" is the phrase to build the contrast on.
%
% AGILE:
%   \parencite[para.~2]{beck2001principles}
%   "Welcome changing requirements, even late in development. Agile processes
%    harness change for the customer's competitive advantage."
%
% OPTIONAL third, and it is a strong one - the traditional authority conceding
% the agile premise and naming Scrum himself:
%   \parencite[p.~379]{kerzner2022}
%   "In traditional project management, we often do not use scope freeze
%    milestones because it is assumed that the scope is well-defined at project
%    initiation. But on other projects, especially in IT, the project may begin
%    based on just an idea and the project's scope must evolve as the project is
%    being implemented. This is quite common with techniques such as Agile and
%    Scrum."
%   !! Do not over-run into the indented block that follows on p. 379 - it is
%      Rudolf Melik (2007), not Kerzner.
%
% OWN VOICE, 2-3 sentences: take ONE named story and show it both ways. Under a
% sequential approach it would have been fixed in a functional specification
% signed off months before any code; in the backlog it was reworded twice and
% re-ordered once after user feedback.
%
% ***** DECIDED: SUB-AXIS 1b IS OUT. *****
% The two Cohn page numbers are unverified, and this is the chapter carrying the
% report's sharpest argument - an unverifiable locator here is the worst possible
% place to take a risk on the one named point-loser. The chapter has five solid
% axes without it, and the outline lists the notation contrast as an "incl.",
% not a requirement.
% TO REINSTATE: verify both pages against a paginated copy of Cohn first. The
% quotes are in the chapter "User Stories Are Not Use Cases" - findable by
% section heading in seconds. If the pages check out, this is a good addition.
%
% OPTIONAL SUB-AXIS (1b) - the notation contrast, use cases/UML vs user stories.
% THE ONLY SOURCE IN THE CORPUS THAT COVERS IT IS COHN:
%   "Use cases are often permanent artifacts that continue to exist as long as
%    the product is under active development or maintenance. Stories, on the
%    other hand, are not intended to outlive the iteration in which they are
%    added to the software."
%   and "use cases are generally written as the result of an analysis activity,
%    while user stories are written as notes that can be used to initiate
%    analysis conversations."
%   !! ******* PAGE NUMBERS NOT VERIFIED (approx. pp. 137-141). *******
%      Check both against a paginated copy of Cohn BEFORE using them, or drop
%      the sub-axis. Kerzner, PMI and Mishra & Alzoubi have ZERO coverage of
%      use cases as a notation, so there is no fallback source.


% --- Axis 2 [CORE - the chosen fifth]: planning horizon ------------------
% TWO CITATIONS.
%
% SEQUENTIAL:
%   \parencite[p.~378]{kerzner2022}
%   "Finally, effective total program planning cannot be accomplished unless all
%    of the necessary information becomes available at project initiation. These
%    information requirements are: the statement of work (SOW); the project
%    specifications; the milestone schedule; the work breakdown structure (WBS)."
%   -> The load-bearing quote for "plan the entire project up front". Then argue,
%      uncited, that none of the four could honestly have been produced at the
%      start of this project.
%
% AGILE:
%   \parencite[ch.~8, pp.~57--62]{kelly2019}
%   "holding rigid priorities, prioritizing into the far future (e.g., more than
%    12 weeks out), and not being prepared to fast-track new valuable work are
%    all behaviors which will cause problems. Such behaviors while appearing to
%    offer certainty and predictability will reduce agility."
%   -> Puts a NUMBER on the horizon, which is what stops this axis being generic.
%
% !! HONEST QUALIFIER, one clause - concede it rather than let it be found:
%    Kerzner says in three separate places that traditional planning is not a
%    single up-front act (rolling wave on p. 305; "project planning is an
%    iterative process and must be performed throughout the life of the project"
%    on p. 367). A reader who checks p. 378 lands three lines below "Stay
%    flexible". Acknowledging this in half a sentence is a Quality point and
%    costs almost nothing; being caught by it is expensive.
%
% OWN VOICE: which parts of this project's plan were genuinely reliable beyond
% the first two sprints, and which were not.


% --- Axis 3 [CORE]: who estimates ----------------------------------------
% TWO CITATIONS - the sharpest inversion in the chapter after axis 5, and the
% one most easily made concrete.
%
% SEQUENTIAL:
%   \parencite[p.~455]{kerzner2022}
%   "Determining the elapsed time between events requires that responsible
%    functional managers evaluate the situation and submit their best estimates."
%   -> In traditional project management the numbers come from MANAGERS, not
%      from the people who will do the work.
%   OPTIONAL reinforcement: \parencite[p.~628]{kerzner2022} "the cost estimates
%   are based on the figures supplied to the project manager by the functional
%   managers."
%   OPTIONAL, sharper still, on who writes the requirements document:
%   \parencite[p.~385]{kerzner2022} "the SOW is prepared by the project office
%   with input from the user groups" - users merely "input".
%
% AGILE:
%   \parencite[p.~11]{scrumguide2020}
%   "The Developers who will be doing the work are responsible for the sizing."
%   -> One sentence, normative, exactly the axis.
%   !! If this sentence was already spent in 3.3 or 3.4, do not spend it a third
%      time - make the agile half in your own words and let the Kerzner citation
%      carry the paragraph.
%
% !! Schwaber (2004) adds NOTHING usable to this axis - his material is
%    narrative and drags in "done" and estimation depth, both out of scope.
%
% OWN VOICE: keep it to the planning consequence, not to estimation technique.
% Effort estimation and planning poker belong to Task 3 - a passing mention at
% most.


% --- Axis 4: NOT USED -----------------------------------------
% Axis 4 (delivery cadence) was considered and dropped; see the decision block at the
% top of this file. The material was removed so it cannot be written by
% mistake. It is preserved in task/Literature review.md and in the
% chapter-4 axis table of PROJECT_FACTS.md if it is ever reinstated.


% --- Axis 5 [CORE - the strongest argument in the report] ----------------
% --- the inverted iron triangle ------------------------------------------
% THREE CITATIONS, and this axis may take slightly more than its share.
%
% SEQUENTIAL, and this is the best material in the corpus because it comes from
% a TRADITIONAL authority rather than an agile advocate:
%   \parencite[p.~78]{kerzner2022}
%   "The waterfall approach, where almost all work is done sequentially, begins
%    with well-defined requirements from which we must determine the budget and
%    schedule to produce the deliverables."
%   immediately followed by
%   "In this case, we may start out with a fixed budget and schedule, and then
%    have to decide how much work can be done within the time and cost
%    constraints. The requirements can evolve over the life of the project."
%   -> ** QUOTE THESE TWO CONSECUTIVELY. ** Together they state the inversion.
%
% ** THE HIGHEST-TRANSFER MOVE AVAILABLE IN THE WHOLE REPORT: **
%   \parencite[p.~609]{kerzner2022} Table 16-1 "Categories of Constraints"
%   enumerates every fixed/variable combination. State that a sequential
%   recipe-app project is Kerzner's A-3 (performance fixed, time and cost float)
%   and a Scrum recipe-app project is his B-1 (time and cost fixed, performance
%   floats). Mapping your own argument onto the source author's taxonomy is
%   exactly what "transfer of theories/models" means.
%   REINFORCEMENT - Kerzner instantiating A-3 himself, so the mapping reads as
%   using his taxonomy rather than inventing a use for it:
%   \parencite[pp.~607--608]{kerzner2022} "situation A-3 portrays most research
%   and development projects. The performance of an R&D project is usually well
%   defined, and it is cost and time that may be allowed to go beyond budget and
%   schedule."
%
% AGILE - only Kelly DERIVES the inversion; Kerzner has 0 occurrences of "iron
% triangle" and PMBOK 0 of "triple constraint":
%   \parencite[ch.~1, pp.~5--10]{kelly2019}
%   The five-step derivation - Cost = People x Time; quality must stay high;
%   time is fixed by the sprint; people are fixed in the short run; therefore
%   "So, people are fixed, costs are fixed, quality needs to be fixed at a high
%   point, long-run deadlines have always been fixed, and short-run deadlines are
%   fixed in agile. That leaves scope."
%   and "If there is one difference between agile working and traditional, it is
%   this: agile constantly flexes scope."
%   -> ONE locator covers the whole derivation, because it is all in chapter 1.
%   !! Include the FIRST step (Cost = People x Time). Without it the chain from
%      "costs are fixed" to "people are fixed" has a hole.
%
% OPTIONAL fourth, PMI's own version, useful because it is not an agile advocacy
% document:
%   \parencite[PMBOK Guide, Section~2.5.7]{pmbok2021}
%   "If the schedule or budget is constrained, the product owner may consider the
%    project done when the highest priority items are delivered."
%
% OWN VOICE, and this is where the axis becomes concrete: name the ONE recipe-app
% feature that gets dropped as the deadline approaches while the core ships. Tie
% it back to the "Could have" band in Table~\ref{tab:initial-backlog} - Cohn's
% definition of "could" ("left out of the release if time runs out", 2004, p. 98)
% was planted in 3.2 exactly so it can be collected here.


% --- Axis 6 [CORE]: feedback and risk ------------------------------------
% --- THIS IS THE ONLY PLACE THE SCRUM CEREMONIES ARE DESCRIBED -----------
% THREE CITATIONS. Refinement and sprint planning were covered in chapter 3 -
% refer back, do not describe them twice.
%
% SEQUENTIAL:
%   \parencite[p.~307]{kerzner2022}
%   Table 8-9, row "Customer feedback": Traditional "Minimal, perhaps only at
%   project termination" / Agile "Throughout the project".
%   OPTIONAL: the phase-gate mechanism, \parencite[pp.~378--379]{kerzner2022} -
%   the end-of-phase meeting with sponsor and senior management, "often called
%   critical design reviews, 'on-off ramps,' and 'gates.'"
%
% AGILE - the Sprint Review, and note WHO attends:
%   \parencite[p.~9]{scrumguide2020}
%   "The Scrum Team presents the results of their work to key stakeholders and
%    progress toward the Product Goal is discussed."
%   -> This is the explicit stakeholder-facing sentence; rest the audience
%      contrast on it.
%
% AGILE - the Retrospective, and note who does NOT attend:
%   \parencite[p.~10]{scrumguide2020}
%   "The Scrum Team inspects how the last Sprint went with regards to
%    individuals, interactions, processes, tools, and their Definition of Done."
%   !! ***** WORD THIS CAREFULLY. ***** The Guide NEVER says stakeholders are
%      excluded. It names only "The Scrum Team". Write "the Guide names only the
%      Scrum Team as participants", NOT "stakeholders are not invited" - the
%      second puts words in the source's mouth.
%
% THE DAILY, one sentence:
%   \parencite[p.~9]{scrumguide2020}
%   "The Daily Scrum is a 15-minute event for the Developers of the Scrum Team."
%
% OWN VOICE: the risk argument - the maximum amount of work that can go in the
% wrong direction is bounded by the review interval. Under a sequential plan,
% the first time a user of the recipe app sees anything is at acceptance.


% --- Axis 7: NOT USED -----------------------------------------
% Axis 7 (communication and transparency) was considered and dropped; see the decision block at the
% top of this file. The material was removed so it cannot be written by
% mistake. It is preserved in task/Literature review.md and in the
% chapter-4 axis table of PROJECT_FACTS.md if it is ever reinstated.


% --- CLOSING: why agile fits THIS product, and where it would not --------
% THREE CITATIONS. This paragraph prepares the balanced ending of chapter 5 -
% do not draw the overall conclusions here, they belong to chapter 5.
%
% WHY AGILE FITS - argue it against a checklist rather than asserting it:
%   \parencite[p.~16]{pmi2017}
%   Projects for which iterative, incremental and agile approaches work well
%   "Require research and development; Have high rates of change; Have unclear or
%    unknown requirements, uncertainty, or risk; or Have a final goal that is
%    hard to describe."
%   -> Apply the four point by point to the recipe app. This is the cleanest way
%      to ARGUE the choice of approach instead of asserting it.
% OWN VOICE: consumer app, evolving user tastes, time to market and first-mover
% advantage.
%
% WHERE SEQUENTIAL WINS - the counter-example. Naming a case where the chosen
% method would be WRONG is a Transfer point.
%   BEST OPTION - PMI's own worked case, so the counter-example is cited rather
%   than asserted:
%   \parencite[p.~138]{pmi2017}
%   "Many lives were potentially at risk, so criticality was very high.
%    Requirements were locked down because changes impacted so many subcontractor
%    organizations."
%   ** AND ITS PAIRED CASE ** - \parencite[p.~136]{pmi2017} describes an online
%   drug store with "extremely volatile requirements with major changes week on
%   week" and very short iterations: a consumer product almost exactly analogous
%   to the recipe app. Citing the PAIR lets the chapter argue both directions
%   from one authority's own examples.
%   ALTERNATIVES:
%     \parencite[p.~1517]{mishra2023} "since safety is a concern, the Department
%     of Defense and the aerospace sector are two industries that would probably
%     choose to employ waterfall over agile." (peer-reviewed backing for the
%     air-traffic-control style example)
%     \parencite[p.~9]{kerzner2022} the Disney case - safety, aesthetic value and
%     quality locked in, all trade-offs made on time, cost and scope.
%   !! Kerzner has NO air-traffic-control or safety-critical example (0 hits) -
%      do not attribute one to him.
%
% BALANCE - one sentence of counter-evidence, and it is worth its space because
% the Quality criterion rewards non-cheerleading:
%   \parencite[p.~1516]{mishra2023}
%   "Agile software development projects take 22.4% longer than anticipated
%    compared to 13.65% longer than the estimated time for the waterfall method"
%   !! Attribute as reported - the figure is Khoza and Marnewick (2020), relayed
%      by Mishra and Alzoubi. This is a narrative review, not an empirical study.
%   !! Take the SCHEDULE figure only. The cost figure from the same passage
%      collides with axis 5 - see trap 1 in the header.
%   !! Be aware the same page credits the same study with 88.2 % vs 47 % success
%      rates in agile's favour. Quoting only the unfavourable half is defensible
%      as balance, but do not pretend the page says only one thing.
%
% OPTIONAL, if the Manifesto's chapter-4 citation has not yet been spent:
%   \parencite[paras.~2--3]{beck2001manifesto}
%   "Responding to change over following a plan" TOGETHER WITH the qualifier
%   "That is, while there is value in the items on the right, we value the items
%    on the left more."
%   -> ** QUOTE THE PAIR. ** The qualifier is what prevents the naive reading
%      "agile means no planning", and handling it explicitly reads as depth. It
%      is also the natural bridge into the hybrid argument in chapter 5.
%      Suggested framing: the Manifesto does not reject planning, it reweights it.
